% ============================================================
% pure 报告 — 规范模板（xelatex + ctex）
% 主题: DDalphaAMG-SM（Schwinger 模型）核心部分穷尽剖析
% 编译: xelatex 两遍; 交付前 Overfull / Float too large 均为 0
% ============================================================
\documentclass[UTF8]{ctexart}
\usepackage[margin=2.5cm]{geometry}
\usepackage{graphicx}
\usepackage{amsmath}
\usepackage{amssymb}
\usepackage{microtype}
\usepackage{listings}
\usepackage{xcolor}
\usepackage{booktabs}
\usepackage{longtable}
\usepackage{xltabular}
\usepackage{tabularx}
\usepackage{hyperref}
\usepackage{fancyhdr}
\usepackage[most]{tcolorbox}
\usepackage{titlesec}

\definecolor{accent}{HTML}{1F4E79}
\definecolor{evidence}{HTML}{1E8449}
\definecolor{reference}{HTML}{8C6D1F}
\definecolor{conclusion}{HTML}{8B1A1A}
\definecolor{codebg}{HTML}{F4F6F7}
\definecolor{algo}{HTML}{E67E22}
\definecolor{phys}{HTML}{6C3483}
\definecolor{map}{HTML}{C2185B}
\definecolor{warn}{HTML}{D35400}
\definecolor{hl}{HTML}{FDF6E3}

\setlength{\emergencystretch}{3em}
\hypersetup{colorlinks=true,linkcolor=accent,urlcolor=map}

\titleformat{\section}{\Large\bfseries\color{accent}}{\thesection}{1em}{}
\titleformat{\subsection}{\large\bfseries\color{phys}}{\thesubsection}{1em}{}
\titleformat{\subsubsection}{\normalsize\bfseries\color{phys!70!black}}{\thesubsubsection}{1em}{}

\newtcolorbox{conclbox}{breakable,colback=conclusion!6,colframe=conclusion,title=结论}
\newtcolorbox{refbox}{breakable,colback=reference!6,colframe=reference,title=参考背景}
\newtcolorbox{warnbox}{breakable,colback=warn!6,colframe=warn,title=局限与风险}
\newtcolorbox{keybox}{breakable,colback=hl,colframe=accent,title=要点}

\lstset{
  basicstyle=\ttfamily\footnotesize,
  backgroundcolor=\color{codebg},
  frame=single,
  language=C++,
  firstnumber=auto,
  keywordstyle=\color{accent}\bfseries,
  commentstyle=\color{evidence!70!black},
  stringstyle=\color{phys},
  breaklines=true,
  breakatwhitespace=false,
  postbreak=\mbox{\textcolor{accent}{$\hookrightarrow$}\space},
  keepspaces=true,
  columns=flexible,
  numbers=left,numberstyle=\tiny\color{phys!60!black},
}

\title{DDalphaAMG-SM（Schwinger 模型）核心部分穷尽剖析报告}
\author{opencode pure}
\date{2026-08-17}

\begin{document}
\maketitle

\begin{abstract}
本报告对 \url{/root/PyQCU/refer/git-rep/DDalphaAMG-SM} 的核心部分进行穷尽剖析。
项目性质判定为\textbf{代码-物理交叉向}：纯 C++/MPI 实现的 2D Schwinger 模型（U(1) 规范）
Wilson 离散化 AMG 求解器。穷尽枚举 9 个核心候选（C1--C9），其中\textbf{深挖 6 / 浅析 2 /
排除 1}：深挖部分含物理图像、公式推导链、代码-物理映射表与证据。独特竞争力：在最小 2D
模型上完整复现 DD$\alpha$AMG 框架——层无关 G 系数算子（粗链接 P$^\dagger$DP）、test-vector
插值、SAP 平滑器、V/K-cycle 与 FGMRES 外层，附完整自校验测试。非独立 git 仓库。
\end{abstract}

\tableofcontents

\section{项目定位与核心部分清单}

\subsection{项目性质判定}
\begin{keybox}
\textbf{性质判定：交叉向。}物理对象明确（2D U(1) Schwinger 模型、Wilson 算子、test vectors、
粗链接），算法代码为核心载体。判定依据：\texttt{src/dirac\_operator.cpp:24-54} 直接实现
Wilson $D$；\texttt{AMG\_Documentation.pdf} 为算法文档；\texttt{README.md:3-13} 明示目的。
\end{keybox}

\subsection{核心部分总清单（穷尽枚举）}
\begin{xltabular}{\textwidth}{lXXX}
\toprule
\textcolor{map}{编号} & 候选 & 三态 & 理由 \\
\midrule
\endhead
C1 & 2D Wilson Dirac 算子 & 深挖 & 求解方程核心，物理对应明确 \\
C2 & test-vector 插值算子 P/P$^\dagger$ & 深挖 & AMG 构造核心 \\
C3 & 层无关 G 系数粗算子 & 深挖 & 本库最独特设计 \\
C4 & SAP 平滑器 & 深挖 & 平滑器核心 \\
C5 & V/K-cycle 循环 & 深挖 & 循环策略核心 \\
C6 & FGMRES 外层与 K-cycle 预条件 & 深挖 & 外层收敛核心 \\
C7 & 粗层 GMRES 与 rank 聚集 & 浅析 & 粗层支持性组件 \\
C8 & 自校验测试族 & 浅析 & 开发期验证，非生产 \\
C9 & 编译期参数固定 & 排除 & 工程约定，无独特性 \\
\bottomrule
\caption{核心部分清单三态（数据来源: 全库索引实测）}
\end{xltabular}

\section{核心部分深度剖析}

\subsection{C1 2D Wilson Dirac 算子}
\textbf{物理图像：}Schwinger 模型（2D U(1)）中夸克为 2 分量旋量，Wilson 离散化：
\begin{equation}\label{eq:wilson2d}
D\psi(x) = (m_0+2)\psi(x) - \frac12\sum_{\mu=0,1}\left[(1-\gamma_\mu)U_\mu(x)\psi(x+\hat\mu) + (1+\gamma_\mu)U_\mu^\dagger(x-\hat\mu)\psi(x-\hat\mu)\right].
\end{equation}
\textbf{代码映射：}
\begin{xltabular}{\textwidth}{lXX}
\toprule
\textcolor{map}{代码符号} & 物理对象 & 含义 \\
\midrule
\endhead
\texttt{U.val[2*n+mu]} & $U_\mu(x)$ & U(1) 复数 link（mu=0 t 向 / mu=1 x 向）\\
\texttt{phi.val} & $\psi(x)$ & 2 分量旋量 \\
\texttt{(m0+2)} & $2+m_0$ & 质量对角项 \\
\texttt{rsign/lsign} & 周期/反周期 & t 向反周期（boundary.h:41-46）\\
\texttt{D\_phi/D\_dagger\_phi} & $D$, $D^\dagger$ & 见 dirac\_operator.cpp:24-85 \\
\bottomrule
\caption{映射表 C1（数据来源: src/*.cpp 实测）}
\end{xltabular}
\textbf{推导链：}Dirac 算子每分量展开见 \texttt{src/dirac\_operator.cpp:37-49}（mu=0/1 两个
分量，含 $\pm i$ 的 $\gamma$ 组合）。\textbf{极限校验：}$m_0=-2$ 时对角项消失；
自由场 $U=1$ 时恢复标准 Wilson 谱。

\subsection{C2 test-vector 插值算子 P/P$^\dagger$}
\textbf{物理图像：}AMG 插值算子 $P$ 由每聚合若干 test vectors 构成，限制为 $P^\dagger$。
\textbf{代码定位：}\texttt{src/level.cpp:92-141}（$P$ 延拓 \texttt{P\_vc}）、
\texttt{:145-189}（$P^\dagger$ 限制 \texttt{Pdagg\_v}）、\texttt{:193-284}
（按块局部 Gram--Schmidt 正交化，Frommer 2014 SIAM 算法）。
\textbf{推导链：}
\begin{equation}\label{eq:p}
\psi_{l} = P_l\,\psi_{l+1},\qquad \eta_{l+1} = P_l^\dagger\,\eta_{l},\qquad P_l^\dagger P_l = I \text{（块内正交）}.
\end{equation}
\textbf{为什么局部正交：}聚合内正交保证 $P$ 列线性无关，避免粗算子病态。

\subsection{C3 层无关 G 系数粗算子}
\textbf{物理图像：}本库最独特设计——Dirac 算子以稀疏 G1/G2/G3 系数（每格点少量复数）
表示，使\textbf{最细层与粗层共用同一 D 应用逻辑}。\textbf{代码定位：}
\texttt{src/level.cpp:286-331}（\texttt{makeDirac} 构造最细层 $G=P=1+\sigma$、$M=1-\sigma$
矩阵）、\texttt{:394-426}（\texttt{D\_operator} 层无关应用）、\texttt{:431-542}
（\texttt{makeCoarseLinks} 粗链接 = 按块聚合 $w^\dagger G w$，区分 A/B/C 系数）。
\textbf{推导链：}粗链接 $G^{l+1} = \sum_{\text{聚合}} w^\dagger G^l w$，粗算子 $D_c = P^\dagger D P$
在系数级别重现。\textbf{为什么这样设计：}2D 模型自由度小，用系数化算子使多层实现统一，
验证性代码可复用同一框架。

\subsection{C4 SAP 平滑器}
\textbf{物理图像：}Schwarz 交替过程：红黑块分解，块内 GMRES 精确求解 $D_B$，块间交替。
\textbf{代码定位：}\texttt{src/sap.cpp:3-41}（\texttt{SchwarzBlocks} 红黑着色）、
\texttt{:48-77}（\texttt{I\_D\_B\_1\_It} 限制→块内求解→延拓）、\texttt{:81-165}
（\texttt{SAP} 主迭代）、\texttt{:169-218}（最细层块内 Dirac \texttt{D\_B}）。
\textbf{推导链：}红黑 Gauss--Seidel 型
\begin{equation}\label{eq:sap}
\psi \leftarrow \psi + \sum_{k\in \text{red}} R_k^T D_{B,k}^{-1} R_k(b - D\psi) \;\rightarrow\; \text{black 同理}.
\end{equation}
\textbf{极限校验：}单块退化为精确求解；迭代收敛时残差单调下降（\texttt{:152-158} 判据）。

\subsection{C5 V/K-cycle 循环}
\textbf{物理图像：}V-cycle：预平滑→残差→限制→递归粗解→延拓→校正→后平滑；
K-cycle：粗层以 FGMRES 包裹。\textbf{代码定位：}\texttt{src/amg.cpp:64-116}（V-cycle）、
\texttt{:119-174}（K-cycle）、\texttt{:177-243}（\texttt{applyMultilevel} 独立迭代入口）。
\textbf{推导链：}V-cycle 一步
\begin{equation}\label{eq:v}
\psi_l \leftarrow S_l^{\nu_2}\big(\psi_l + P_l\, V_{l+1}\, P_l^\dagger(\eta_l - D_l\psi_l)\big),\quad \text{经 } S_l^{\nu_1} \text{ 预平滑}.
\end{equation}
\textbf{为什么 K-cycle：}以 FGMRES 重组合粗层（\texttt{\nolinkurl{src/amg.cpp:154}}
\nolinkurl{fgmres_k_cycle_l[l+1]->fgmres(...)}），提升收敛率。

\subsection{C6 FGMRES 外层与 K-cycle 预条件}
\textbf{物理图像：}外层 FGMRES 以 AMG（V/K-cycle）为右预条件子，灵活内积保存预条件子。
\textbf{代码定位：}\texttt{src/fgmres.cpp:16-92}（Arnoldi）、\texttt{:99-115}
（Givens 旋转）、\texttt{:118-128}（上三角回代）；\texttt{include/amg.h:147-270}
（FGMRES\_AMG\_v/k\_cycle 类族）。
\textbf{为什么 FGMRES 而非 GMRES：}变预条件（如 K-cycle 内层收敛变化）下 FGMRES 保留
收敛性，适合 AMG 预条件组合。

\section{独特性与竞争力评估}
\begin{xltabular}{\textwidth}{lXX}
\toprule
\textcolor{algo}{核心} & 独特竞争力 & 对比朴素实现 \\
\midrule
\endhead
层无关 G 系数算子 & 多层共用同一 D 逻辑 & 每层独立写算子 \\
test-vector 插值 & 近零模捕捉 + 局部正交 & 静态粗化 \\
SAP + V/K-cycle & 完整 AMG 框架于 2D 复现 & 单层 Krylov \\
FGMRES 预条件 & 灵活内积容变预条件 & GMRES 固定预条件 \\
自校验测试族 & Dc=P$^\dagger$DP、P$^\dagger$P=I 验证 & 无验证 \\
\bottomrule
\caption{独特性评估（数据来源: 源码实测）}
\end{xltabular}

\section{关键片段与公式}
\begin{lstlisting}[language=C++]
void AlgebraicMG::v_cycle(const int& l, const spinor& eta_l, spinor& psi_l){
	if (l == LevelV::maxLevel && levels[l]->ranks_comm != MPI_COMM_NULL){
		levels[l]->gmres_l->fgmres(eta_l, eta_l, psi_l, false);
	} else {
		if (nu1 > 0) levels[l]->sap_l->SAP(eta_l,psi_l,nu1,tol,false);
		levels[l]->D_operator(psi_l,Dpsi);
		r_l = eta_l - Dpsi;
		levels[l]->Pdagg_v(r_l,eta_l_1);
		if (levels[l+1]->ranks_comm != MPI_COMM_NULL)
			v_cycle(l+1,eta_l_1,psi_l_1);
		levels[l]->P_vc(psi_l_1,P_psi);
		psi_l += P_psi;
		if (AMGV::nu2 > 0) levels[l]->sap_l->SAP(eta_l,psi_l,nu2,tol,false);
	}
}
\end{lstlisting}
参考源：\texttt{src/amg.cpp:64-116}（V-cycle，注释为简化说明）

核心公式汇总：
\begin{align}
D &= (m_0+2)\mathbb{1} - \tfrac12\sum_\mu[(1-\gamma_\mu)U_\mu\delta_{+}+(1+\gamma_\mu)U_\mu^\dagger\delta_{-}], \label{eq:f1}\\
D_c &= P^\dagger D P \quad(\text{系数级 G1/G2/G3}), \label{eq:f2}\\
P_l^\dagger P_l &= I \quad(\text{块内正交}). \label{eq:f3}
\end{align}

\section{参考源清单}
{\small
\begin{xltabular}{\textwidth}{XXX}
\toprule
\textcolor{evidence}{参考源} & 类型 & 内容/作用 \\
\midrule
\endhead
\texttt{\nolinkurl{src/dirac_operator.cpp:24-92}} & 仓库 & $D$/$D^\dagger$/$D^\dagger D$ \\
\texttt{\nolinkurl{src/level.cpp:92-189}} & 仓库 & P/P$^\dagger$ 转移 \\
\texttt{\nolinkurl{src/level.cpp:193-284}} & 仓库 & 局部 Gram--Schmidt \\
\texttt{\nolinkurl{src/level.cpp:286-331}} & 仓库 & makeDirac 最细层 G 系数 \\
\texttt{\nolinkurl{src/level.cpp:394-426}} & 仓库 & 层无关 D 应用 \\
\texttt{\nolinkurl{src/level.cpp:431-542}} & 仓库 & 粗链接 P$^\dagger$DP \\
\texttt{\nolinkurl{src/sap.cpp:3-41}} & 仓库 & Schwarz 块/红黑着色 \\
\texttt{\nolinkurl{src/sap.cpp:48-77}} & 仓库 & 限制-块内 GMRES-延拓 \\
\texttt{\nolinkurl{src/sap.cpp:81-165}} & 仓库 & SAP 主迭代 \\
\texttt{\nolinkurl{src/amg.cpp:64-116}} & 仓库 & V-cycle \\
\texttt{\nolinkurl{src/amg.cpp:119-174}} & 仓库 & K-cycle \\
\texttt{\nolinkurl{src/amg.cpp:177-243}} & 仓库 & applyMultilevel \\
\texttt{\nolinkurl{src/fgmres.cpp:16-128}} & 仓库 & FGMRES 骨架 \\
\texttt{\nolinkurl{include/amg.h:147-270}} & 仓库 & FGMRES\_AMG 类族 \\
\texttt{\nolinkurl{include/boundary.h:41-46}} & 仓库 & t 向反周期 \\
\texttt{\nolinkurl{src/io.cpp:32-46}} & 仓库 & 二进制配置读取 \\
\texttt{\nolinkurl{README.md:3-13}} & 仓库 & 项目定位 \\
\bottomrule
\end{xltabular}
}

\section{结论}
\begin{conclbox}
DDalphaAMG-SM 的核心竞争力是\textbf{层无关 G 系数粗算子设计}：在最小 2D U(1) 模型上以系数化
算子统一最细/粗层 Dirac 应用，配合 test-vector 插值、SAP、V/K-cycle 与 FGMRES 预条件，完整
复现 DD$\alpha$AMG 框架并附详尽自校验。穷尽枚举 9 个候选，深挖 6 项均有物理图像与代码-物理
映射证据。
\end{conclbox}
\begin{warnbox}
\textcolor{warn}{局限：}未实测运行（无构建产物）；rank 聚集硬编码 4（variables.h:85-87）；
数值参数编译期固定；非独立 git 仓库。
\end{warnbox}

\end{document}